% =====================================================================
%  Thème 3 — Colinéarité et parallélisme — Niveau DIFFICILE — Corrigé
% =====================================================================
\documentclass[11pt,a4paper]{article}
\usepackage[utf8]{inputenc}
\usepackage[T1]{fontenc}
\usepackage{lmodern}
\usepackage[french]{babel}
\usepackage{amsmath,amssymb,mathtools}
\usepackage{xcolor}
\usepackage{enumitem}
\usepackage{fancyhdr}
\usepackage[a4paper,margin=2.1cm,headheight=26pt,headsep=10pt]{geometry}

\definecolor{primary}{RGB}{222,70,80}
\definecolor{primarydark}{RGB}{150,30,42}
\definecolor{excol}{RGB}{0,135,90}

\pagestyle{fancy}\fancyhf{}
\fancyhead[L]{\small\textcolor{primarydark}{\textbf{Fiche 2} — Calcul vectoriel}}
\fancyhead[R]{\small\textcolor{primarydark}{\textsc{Thème 3 — Difficile — Corrigé}}}
\fancyfoot[C]{\small\thepage}
\renewcommand{\headrulewidth}{0.8pt}

\newcommand{\vc}[1]{\overrightarrow{#1}}
\providecommand{\norm}[1]{\left\lVert #1\right\rVert}
\newcommand{\cor}[1]{\medskip\noindent{\color{primarydark}\bfseries\large Corrigé #1.}\par\smallskip}
\newcommand{\rep}{\textcolor{excol}{$\blacktriangleright$}\ }

\begin{document}
\begin{center}
{\LARGE\bfseries\color{primarydark}Colinéarité et parallélisme — Corrigé Difficile}\\[2pt]
{\color{primary}\rule{0.5\linewidth}{1.1pt}}
\end{center}
\vspace{2mm}

\cor{1}
\begin{enumerate}[label=\textbf{\alph*)},leftmargin=2em,itemsep=3pt]
  \item Déterminant de $\vec u\,(\sqrt3;1)$ et $\vec v\,(3;\sqrt3)$ :
        $\sqrt3\times\sqrt3-3\times1=3-3=0$. Ils sont colinéaires. De plus
        $\vec v=(3;\sqrt3)=\sqrt3\,(\sqrt3;1)=\sqrt3\,\vec u$, donc $k=\sqrt3$.
  \item Déterminant de $\vec u\,(\sqrt3;1)$ et $\vec w\,(2;2\sqrt3)$ :
        $\sqrt3\times2\sqrt3-2\times1=2\times3-2=6-2=4\neq0$. Donc $\vec w$ n'est \textbf{pas}
        colinéaire à $\vec u$.
\end{enumerate}
\rep Avec des radicaux, on garde $\sqrt3\times\sqrt3=3$ : le déterminant redevient un nombre simple.

\cor{2}
\begin{enumerate}[label=\textbf{\alph*)},leftmargin=2em,itemsep=3pt]
  \item Déterminant $=m\times m-8\times2=m^2-16$.
  \item Colinéaires $\iff m^2-16=0 \iff m^2=16 \iff m=4 \text{ ou } m=-4$.
  \item Pour $m=4$ : $\vec u\,(4;2)$, $\vec v\,(8;4)=2\,\vec u$, donc $k=2$.\\
        Pour $m=-4$ : $\vec u\,(-4;2)$, $\vec v\,(8;-4)=-2\,\vec u$, donc $k=-2$.
\end{enumerate}
\rep Une condition de colinéarité donne ici une équation du \emph{second degré} : deux valeurs de $m$.

\cor{3}
\begin{enumerate}[label=\textbf{\alph*)},leftmargin=2em,itemsep=3pt]
  \item $\vc{AB}\,(3;2)$ et $\vc{AC}\,(9;6)$. Déterminant $=3\times6-9\times2=18-18=0$ :
        $A$, $B$, $C$ sont alignés.
  \item $D\,(7;y_D)$ sur la droite $\iff \vc{AD}$ colinéaire à $\vc{AB}$. Or $\vc{AD}\,(6;\,y_D-1)$ et
        $\vc{AB}\,(3;2)$ :
        \[6\times2-3\times(y_D-1)=12-3y_D+3=15-3y_D=0 \Rightarrow y_D=5.\]
        L'antenne est en $D\,(7;5)$.
  \item $\vc{BD}\,(7-4;5-3)=(3;2)=\vc{AB}$ : $\vc{BD}$ est bien colinéaire à $\vc{AB}$ (ici même
        vecteur), $D$ est donc sur la droite $(AB)$.
\end{enumerate}
\rep « Être sur la droite $(AB)$ » se traduit par « $\vc{AD}$ colinéaire à $\vc{AB}$ », c'est-à-dire
déterminant nul.

\cor{4}
\begin{enumerate}[label=\textbf{\Alph*.},leftmargin=2.2em,itemsep=3pt]
  \item \textbf{Vrai}. $\vc{AB}\,(1;2)$ et $\vc{AC}\,(-2;-4)=-2\,(1;2)=-2\,\vc{AB}$ : colinéaires, donc
        $A$, $B$, $C$ alignés.
  \item \textbf{Vrai}. $\vc{AB}\,(1;2)$ est un vecteur directeur de $(AB)$, et $\vec u\,(1;2)$ lui est
        égal.
  \item \textbf{Vrai}. Pour $d:\ 4x-2y+5=0$, $a=4$, $b=-2$, un vecteur directeur est
        $(-b;a)=(2;4)=2\,(1;2)$, colinéaire à $\vc{AB}$ : $(AB)\parallel d$.
  \item \textbf{Faux}. La droite $(AB)$ a pour équation $y=2x+3$, soit $2x-y+3=0$, ou encore
        $4x-2y+6=0$. Elle diffère de $d$ ($4x-2y+5=0$) par la constante ($6\neq5$) : les droites sont
        \emph{strictement parallèles}, pas confondues.
\end{enumerate}
\rep Parallèles \emph{et} un point commun $\Rightarrow$ confondues. Ici aucun point commun : parallèles
strictes.

\end{document}
